\SourceLeaf{67}
$\delta$) Setzt man also

\[
\left(\frac{1-\BareRoot{(1-nn)}}{n}\right) = a,
\]

\[
\begin{aligned}
\log(1+n\cos\varphi) & = -\log\frac{2a}{n} + \frac{2}{1}a\cos\varphi - \frac{2}{2}a^2\cos2\varphi +\\
&{} + \frac{2}{3}a^3\cos3\varphi + .\,.\,.\,. + (-1)^{h+1}\frac{2}{h}a^h\cos h\varphi + .\,.\,.\,.
\end{aligned}
\]

\textls[140]{Anmerkung.} Diese Formel ist für die Bestimmung der Logarithmen der trigonometrischen
Funktionen von größtem Nutzen.

Denn setzt man $n=1$, so wird $a=1$, und man erhält

\[
\log(1+\cos\varphi) = -\log 2 + \frac{2}{1}\cos\varphi - \frac{2}{2}\cos2\varphi + \frac{2}{3}\cos3\varphi - \frac{2}{4}\cos4\varphi + .\,.\,.,
\]

\noindent Diese Reihe liefert, da $1+\cos\varphi = 2\cos\frac{1}{2}\varphi^2$,
auch

\[
\log\cos\tfrac{1}{2}\varphi = -\log 2 + \cos\varphi - \tfrac{1}{2}\cos2\varphi + \tfrac{1}{3}\cos3\varphi - \tfrac{1}{4}\cos4\varphi + .\,.\,.
\]

\noindent und somit durch Einsetzen von $2\varphi$ anstelle von $\varphi$

\[
\log\cos\varphi = -\log 2 + \cos2\varphi - \tfrac{1}{2}\cos4\varphi + \tfrac{1}{3}\cos6\varphi - \tfrac{1}{4}\cos8\varphi + .\,.\,.\,.
\]

\noindent Ebenso erhält man, wenn man $n=-1$ setzt, sodass $a=-1$ wird,

\[
\log(1-\cos\varphi) = -\log 2 - \frac{2}{1}\cos\varphi - \frac{2}{2}\cos2\varphi - \frac{2}{3}\cos3\varphi - \frac{2}{4}\cos4\varphi - .\,.\,.\,.
\]

\noindent und somit, da $1-\cos\varphi = 2\sin\tfrac{1}{2}\varphi^2$,

\[
\log\sin\tfrac{1}{2}\varphi = -\log 2 - \cos\varphi - \tfrac{1}{2}\cos2\varphi - \tfrac{1}{3}\cos3\varphi - \tfrac{1}{4}\cos4\varphi - .\,.\,.\,.
\]

\noindent oder, indem man wie zuvor $2\varphi$ anstelle von $\varphi$ einsetzt,

\[
\log\sin\varphi = -\log 2 - \cos2\varphi - \tfrac{1}{2}\cos4\varphi - \tfrac{1}{3}\cos6\varphi - \tfrac{1}{4}\cos8\varphi - .\,.\,.\,.
\]

\noindent Hieraus folgt wegen $\log\operatorname{tang}\varphi = \log\sin\varphi - \log\cos\varphi$

\[
\log\operatorname{tang}\varphi = -2\left(\cos2\varphi + \tfrac{1}{3}\cos6\varphi + \tfrac{1}{5}\cos10\varphi + \tfrac{1}{7}\cos14\varphi + .\,.\,.\right)
\]

\noindent und auf demselben Weg erhält man die Logarithmen der übrigen trigonometrischen Funktionen, wenn auch durch
langsam konvergierende Reihen.

$\varepsilon$) Da

\[
\log(1+n\cos\varphi) = -\log\frac{2a}{n} + \frac{2}{1}a\cos\varphi - \frac{2}{2}a^2\cos2\varphi + \frac{2}{3}a^3\cos3\varphi - .\,.\,.\,.\,.,
\]

\noindent ist die Integration des Ausdrucks $\log(1+n\cos\varphi).\,d\varphi$ unmittelbar möglich und ergibt

\[
\int d\varphi\,\log(1+n\cos\varphi) = -\left(\log\frac{2a}{n}\right)\varphi + \frac{2}{1}a\sin\varphi - \frac{2}{4}a^2\sin2\varphi + \frac{2}{9}a^3\sin3\varphi - .\,.\,. 
\]

